\documentclass[a4paper,11pt]{article}
\usepackage{lmodern}
\usepackage[T1]{fontenc}
\usepackage[utf8]{inputenc}
\usepackage{amsmath,amssymb,amsthm}
\usepackage{longtable,booktabs,array,calc}
\usepackage{hyperref}
\usepackage[margin=25mm]{geometry}

\newtheorem{theorem}{Theorem}[section]
\newtheorem{proposition}[theorem]{Proposition}
\newtheorem{lemma}[theorem]{Lemma}
\newtheorem{corollary}[theorem]{Corollary}
\newtheorem{definition}[theorem]{Definition}
\newtheorem{remark}[theorem]{Remark}
\newtheorem{observation}[theorem]{Observation}

\providecommand{\tightlist}{\setlength{\itemsep}{0pt}\setlength{\parskip}{0pt}}
\newcounter{none}

\begin{document}

\section{Combinatorial Properties of the Configuration Structure of the
6-Dimensional
Hyperrectangle}\label{combinatorial-properties-of-the-configuration-structure-of-the-6-dimensional-hyperrectangle}

\subsection{------ Spin Classification, Configuration Enumeration, and
Signed Area
------}\label{spin-classification-configuration-enumeration-and-signed-area}

\textbf{Author}: Noriaki Kihara (WF System Co., Ltd.) \textbf{Date}:
April 2026 \textbf{DOI}: 10.5281/zenodo.19646652

\begin{center}\rule{0.5\linewidth}{0.5pt}\end{center}

\subsection{0. Purpose of This Paper}\label{purpose-of-this-paper}

This paper defines a 6-dimensional structure obtained by adding a
discrete label axis \(Q\) to the 5-dimensional hyperrectangle, and
derives the combinatorial properties arising from its configuration
structure.

The main body (§1--§9) states only the mathematical structure of the
model and the consequences derived purely from combinatorics. The
correspondence with the Standard Model is presented in §10 as ``an
example of interpretation,'' and does not constitute a claim of this
paper. It is merely an example showing how the configuration structure
of the model might be mapped to the particle classification of the
Standard Model.

\textbf{On the external assumptions of this model}: This paper takes the
5-dimensional hyperrectangle as its starting point. Why 5 dimensions,
and why 3 of the 5 axes are later treated as ``spatial axes,'' are
\textbf{external assumptions of the model} and are not justified within
the model itself. This dimensional setting is adopted as a
correspondence to the empirical fact that the observed spacetime
consists of 3 spatial dimensions + 1 temporal dimension.

\begin{center}\rule{0.5\linewidth}{0.5pt}\end{center}

\subsection{1. Definition of the 6-Dimensional
Structure}\label{definition-of-the-6-dimensional-structure}

\subsubsection{1.1 Definition}\label{definition}

The 6-dimensional structure treated in this paper is defined as follows.

\begin{itemize}
\tightlist
\item
  \textbf{Geometric structure}: The combinatorial structure of the
  5-dimensional hyperrectangle (vertices, edges, faces, 3-cells,
  4-facets, and 5-cell body)
\item
  \textbf{Label axis}: A discrete label axis \(Q \in \{0, r, g, b\}\)
  with 4 values
\end{itemize}

The vertex coordinates of the 5-dimensional hyperrectangle are denoted
\((x_1, x_2, x_3, x_4, x_5)\), where each coordinate \textbf{takes any
real value}. When a specific proposition imposes conditions on
coordinate values within the configuration structure, such conditions
are declared locally within that proposition.

The 6th axis \(Q\) is not a geometric coordinate but a discrete label
value associated with each state. Its range is the 4 values
\(\{0, r, g, b\}\). The addition of the \(Q\) axis does not change the
combinatorial structure (number of vertices, edges, faces, etc.) of the
hyperrectangle.

However, in the classification of bivectors that determine spin, \(Q\)
participates in the enumeration of orientations as one axis on equal
footing with the other 5 axes. That is, combinations such as
\(e_t \wedge e_Q\) and \(e_R \wedge e_Q\) are abstract bivector pairs
defined as axis combinations, not geometric 2-faces.

\subsubsection{\texorpdfstring{1.2 Element Counts of the
\(k\)-Dimensional
Hyperrectangle}{1.2 Element Counts of the k-Dimensional Hyperrectangle}}\label{element-counts-of-the-k-dimensional-hyperrectangle}

Let \(f_j(k)\) denote the number of \(j\)-dimensional faces
(\(j\)-faces) of the \(k\)-dimensional hyperrectangle. We write out the
element counts from low dimensions so that the reader can verify the
general formula.

\textbf{Vertex count \(f_0(k)\)}: The combinatorial count distinguishing
the 2 states at both ends of each axis

{\def\LTcaptype{none} % do not increment counter
\begin{longtable}[]{@{}ccccccc@{}}
\toprule\noalign{}
\(k\) & 0 & 1 & 2 & 3 & 4 & 5 \\
\midrule\noalign{}
\endhead
\bottomrule\noalign{}
\endlastfoot
\(f_0(k)\) & 1 & 2 & 4 & 8 & 16 & 32 \\
\end{longtable}
}

\[f_0(k) = 2^k\]

\textbf{Edge count \(f_1(k)\)}: Extending along 1 axis, determined by
the positions at both ends of the remaining \((k-1)\) axes

{\def\LTcaptype{none} % do not increment counter
\begin{longtable}[]{@{}ccccccc@{}}
\toprule\noalign{}
\(k\) & 0 & 1 & 2 & 3 & 4 & 5 \\
\midrule\noalign{}
\endhead
\bottomrule\noalign{}
\endlastfoot
\(f_1(k)\) & 0 & 1 & 4 & 12 & 32 & 80 \\
\end{longtable}
}

\[f_1(k) = k \cdot 2^{k-1}\]

Verification: For \(k=3\), \(3 \cdot 4 = 12\) (12 edges of a cube); for
\(k=5\), \(5 \cdot 16 = 80\).

\textbf{Square face count \(f_2(k)\)}: Faces spanned by 2 axes

{\def\LTcaptype{none} % do not increment counter
\begin{longtable}[]{@{}ccccccc@{}}
\toprule\noalign{}
\(k\) & 0 & 1 & 2 & 3 & 4 & 5 \\
\midrule\noalign{}
\endhead
\bottomrule\noalign{}
\endlastfoot
\(f_2(k)\) & 0 & 0 & 1 & 6 & 24 & 80 \\
\end{longtable}
}

\[f_2(k) = \binom{k}{2} \cdot 2^{k-2}\]

Verification: For \(k=3\), \(3 \cdot 2 = 6\) (6 faces of a cube); for
\(k=5\), \(10 \cdot 8 = 80\).

\textbf{Cubic cell (3-cell) count \(f_3(k)\)}: Cubic cells spanned by 3
axes

{\def\LTcaptype{none} % do not increment counter
\begin{longtable}[]{@{}ccccccc@{}}
\toprule\noalign{}
\(k\) & 0 & 1 & 2 & 3 & 4 & 5 \\
\midrule\noalign{}
\endhead
\bottomrule\noalign{}
\endlastfoot
\(f_3(k)\) & 0 & 0 & 0 & 1 & 8 & 40 \\
\end{longtable}
}

\[f_3(k) = \binom{k}{3} \cdot 2^{k-3}\]

Verification: For \(k=4\), \(4 \cdot 2 = 8\); for \(k=5\),
\(10 \cdot 4 = 40\).

\textbf{4-facet count \(f_4(k)\)}: 4-dimensional cells spanned by 4 axes

{\def\LTcaptype{none} % do not increment counter
\begin{longtable}[]{@{}ccccccc@{}}
\toprule\noalign{}
\(k\) & 0 & 1 & 2 & 3 & 4 & 5 \\
\midrule\noalign{}
\endhead
\bottomrule\noalign{}
\endlastfoot
\(f_4(k)\) & 0 & 0 & 0 & 0 & 1 & 10 \\
\end{longtable}
}

\[f_4(k) = \binom{k}{4} \cdot 2^{k-4}\]

Verification: For \(k=5\), \(5 \cdot 2 = 10\).

\textbf{5-cell body count \(f_5(k)\)}: For \(k=5\)

\[f_5(5) = \binom{5}{5} \cdot 2^0 = 1\]

\subsubsection{General Formula}\label{general-formula}

Summarizing the above, the number of \(j\)-dimensional faces of the
\(k\)-dimensional hyperrectangle is given by

\[f_j(k) = \binom{k}{j} \cdot 2^{k-j} \qquad (0 \leq j \leq k)\]

This corresponds to choosing \(j\) axes out of \(k\) to span the
\(j\)-face (\(\binom{k}{j}\) ways) and specifying which end of each of
the remaining \((k-j)\) axes the face is located at (\(2^{k-j}\) ways).

\textbf{Verification (Euler characteristic)}:
\(\sum_{j=0}^{k} (-1)^j f_j(k) = (2-1)^k = 1\), which is consistent with
the hyperrectangle being a contractible space.

\subsubsection{\texorpdfstring{Values Used in This Paper
(\(k=5\))}{Values Used in This Paper (k=5)}}\label{values-used-in-this-paper-k5}

{\def\LTcaptype{none} % do not increment counter
\begin{longtable}[]{@{}clc@{}}
\toprule\noalign{}
\(j\) & Element & \(f_j(5)\) \\
\midrule\noalign{}
\endhead
\bottomrule\noalign{}
\endlastfoot
0 & Vertices & 32 \\
1 & Edges & 80 \\
2 & Square faces & 80 \\
3 & Cubic cells (3-cells) & 40 \\
4 & 4-facets & 10 \\
5 & 5-cell body & 1 \\
\end{longtable}
}

\subsubsection{1.3 Axis Symmetry and Definition of Sign
Configurations}\label{axis-symmetry-and-definition-of-sign-configurations}

\paragraph{1.3.1 Complete Symmetry of 5
Axes}\label{complete-symmetry-of-5-axes}

\textbf{Proposition 1.1 (Complete symmetry of 5 axes)}. The 5 geometric
axes of the 5-dimensional hyperrectangle are all equivalent in their
combinatorial structure. Any permutation of axes gives an automorphism
of the hyperrectangle, and the 5 axes are indistinguishable with respect
to the element counts \(f_j(k)\) of §1.2 and subsequent geometric
propositions.

\textbf{Proof}. The automorphism group of the combinatorial structure of
the \(k\)-dimensional hyperrectangle is the hyperoctahedral group
\(B_k\), which includes all axis permutations. Therefore, the 5 axes are
completely symmetric combinatorially. \(\square\)

\paragraph{1.3.2 Axis Classification}\label{axis-classification}

\textbf{Definition 1.2 (Axis notation labels)}. The following notation
labels are assigned to the 5 axes \(x_1, x_2, x_3, x_4, x_5\).

\begin{itemize}
\tightlist
\item
  \(x_1, x_2, x_3\) are denoted \(x, y, z\)
\item
  \(x_4\) is denoted \(t\)
\item
  \(x_5\) is denoted \(R\)
\end{itemize}

By Proposition 1.1, the 5 axes are equivalent in combinatorial
structure, so the assignment of labels to axes is merely a definitional
choice. These labels are used as notational prerequisites for the
definition of sign configurations (§1.3.3) and the definition of spin
(§2).

\textbf{\(x, y, z, t, R\) are merely notational labels and do not imply
geometric distinction among the axes.} The assignment of physical
interpretations (space, time, curvature, etc.) to each axis is performed
in §10 (an example of interpretation), which is an interpretation from
outside the model and is independent of the geometric structure of §1.

\paragraph{1.3.3 Rigorous Definition of Sign
Configurations}\label{rigorous-definition-of-sign-configurations}

Each component \(x_i\) (\(i = 1, \ldots, 5\)) of a geometric axis is a
symbol representing the \textbf{range constraint} that the coordinate
value along that axis must satisfy. This paper uses the following 5
types of symbols.

{\def\LTcaptype{none} % do not increment counter
\begin{longtable}[]{@{}cll@{}}
\toprule\noalign{}
Symbol & Meaning & Range of coordinate values \\
\midrule\noalign{}
\endhead
\bottomrule\noalign{}
\endlastfoot
\(+\) & Positive & \(x_i > 0\) \\
\(-\) & Negative & \(x_i < 0\) \\
\(0\) & Origin & \(x_i = 0\) \\
\(\pm\) & Non-zero & \(x_i \neq 0\) (positive or negative) \\
\(*\) & Arbitrary & \(x_i\) is any real number \\
\end{longtable}
}

For the \(Q\) axis, similar symbols representing the range are used.

{\def\LTcaptype{none} % do not increment counter
\begin{longtable}[]{@{}
  >{\centering\arraybackslash}p{(\linewidth - 4\tabcolsep) * \real{0.1786}}
  >{\raggedright\arraybackslash}p{(\linewidth - 4\tabcolsep) * \real{0.3214}}
  >{\raggedright\arraybackslash}p{(\linewidth - 4\tabcolsep) * \real{0.5000}}@{}}
\toprule\noalign{}
\begin{minipage}[b]{\linewidth}\centering
Symbol
\end{minipage} & \begin{minipage}[b]{\linewidth}\raggedright
Meaning
\end{minipage} & \begin{minipage}[b]{\linewidth}\raggedright
Value of \(Q\)
\end{minipage} \\
\midrule\noalign{}
\endhead
\bottomrule\noalign{}
\endlastfoot
\(0\) & Color-neutral & \(Q = 0\) \\
\(r, g, b\) & Specified color & \(Q\) is the specified color label \\
\(r/g/b\) & Any of the 3 colors & \(Q \in \{r, g, b\}\) \\
\(\bar{Q}\) & Anti-color &
\(r \leftrightarrow \bar{r}, g \leftrightarrow \bar{g}, b \leftrightarrow \bar{b}, 0 \leftrightarrow 0\) \\
\(*\) & Arbitrary & \(Q \in \{0, r, g, b\}\) \\
\end{longtable}
}

\textbf{Remark}. Sign configurations are not coordinate values
themselves but symbols specifying \textbf{conditions} that coordinate
values must satisfy. As stated in §1.1, coordinate values in this model
can take any real value, but in specific propositions and configuration
structures, the above range constraints are imposed on the coordinate
values of the relevant axes.

\textbf{Example}. The configuration \((+, +, 0, 0, 0, 0)\) refers to the
set of states satisfying
\(x_1 > 0, x_2 > 0, x_3 = 0, x_4 = 0, x_5 = 0, Q = 0\).

\begin{center}\rule{0.5\linewidth}{0.5pt}\end{center}

\subsection{2. Definition of Spin}\label{definition-of-spin}

\textbf{Definition 2.1 (Spin)}. For a configuration
\((x, y, z, t, R, Q)\), let \(n\) denote the number of axes among
\(t, R, Q\) whose sign configuration is not \(0\). Spin \(s\) is defined
as follows.

\begin{itemize}
\tightlist
\item
  When 2 of the \(xyz\) axes are non-zero and 1 is \(0\):
  \(s = n + 1/2\) (fermion)
\item
  Otherwise: \(s = n\) (boson)
\end{itemize}

Since \(t, R, Q\) comprise 3 axes, the possible values of \(n\) are the
4 cases \(0, 1, 2, 3\). Therefore, the possible values of spin \(s\) in
this model are the following 8 types.

{\def\LTcaptype{none} % do not increment counter
\begin{longtable}[]{@{}ccc@{}}
\toprule\noalign{}
\(n\) & Boson \(s = n\) & Fermion \(s = n + 1/2\) \\
\midrule\noalign{}
\endhead
\bottomrule\noalign{}
\endlastfoot
0 & \(s = 0\) & \(s = 1/2\) \\
1 & \(s = 1\) & \(s = 3/2\) \\
2 & \(s = 2\) & \(s = 5/2\) \\
3 & \(s = 3\) & \(s = 7/2\) \\
\end{longtable}
}

\textbf{Remark}. This definition is a combinatorial definition internal
to the model. The terms ``boson'' and ``fermion'' here are terminology
within the model; correspondence to physical spin angular momentum,
statistics, \((-1)^{2s}\) transformation properties, etc. is treated in
§10 (an example of interpretation).

\begin{center}\rule{0.5\linewidth}{0.5pt}\end{center}

\subsection{\texorpdfstring{3. Enumeration of \(Q\)-Label
Transitions}{3. Enumeration of Q-Label Transitions}}\label{enumeration-of-q-label-transitions}

Consider transitions among the 3 non-zero values \(\{r, g, b\}\) of the
\(Q\) axis. A transition is described by the pair
\((Q_\text{initial}, Q_\text{final})\), and the total number of
transitions from 3 values to 3 values is

\[3 \times 3 = 9\]

Among these, one linear combination of the diagonal components (the
state proportional to \(r\bar{r} + g\bar{g} + b\bar{b}\)) has the same
action on all color labels and does not mediate label transitions.
Excluding this 1 state,

\[3 \times 3 - 1 = 8\]

is the number of independent transition states. Physical naming is given
in §10.

\begin{center}\rule{0.5\linewidth}{0.5pt}\end{center}

\subsection{4. Enumeration of Spin-1
Configurations}\label{enumeration-of-spin-1-configurations}

Spin-1 configurations have \(n = 1\) by Definition 2.1, meaning exactly
1 of the \(t, R, Q\) axes has a non-zero sign configuration, while all
\(xyz\) axes are \(0\).

Depending on which of \(t, R, Q\) is non-zero, the number of
configuration states is classified as follows.

{\def\LTcaptype{none} % do not increment counter
\begin{longtable}[]{@{}
  >{\raggedright\arraybackslash}p{(\linewidth - 4\tabcolsep) * \real{0.2692}}
  >{\raggedright\arraybackslash}p{(\linewidth - 4\tabcolsep) * \real{0.3846}}
  >{\centering\arraybackslash}p{(\linewidth - 4\tabcolsep) * \real{0.3462}}@{}}
\toprule\noalign{}
\begin{minipage}[b]{\linewidth}\raggedright
Non-zero axis
\end{minipage} & \begin{minipage}[b]{\linewidth}\raggedright
Configuration form
\end{minipage} & \begin{minipage}[b]{\linewidth}\centering
Number of states
\end{minipage} \\
\midrule\noalign{}
\endhead
\bottomrule\noalign{}
\endlastfoot
\(t\) & \((0, 0, 0, \pm, 0, 0)\) & 2 (2 choices of sign \(\pm\)) \\
\(R\) & \((0, 0, 0, 0, \pm, 0)\) & 2 (2 choices of sign \(\pm\)) \\
\(Q\) & \((0, 0, 0, 0, 0, Q_\text{initial} \to Q_\text{final})\) & 9 (by
§3) \\
\end{longtable}
}

\textbf{Total}: \(2 + 2 + 9 = 13\) states (for the interpretation
excluding 1 state, see §10)

\begin{center}\rule{0.5\linewidth}{0.5pt}\end{center}

\subsection{5. Enumeration of Spin-2
Configurations}\label{enumeration-of-spin-2-configurations}

Spin-2 configurations have \(n = 2\) by Definition 2.1, meaning exactly
2 of the \(t, R, Q\) axes have non-zero sign configurations, while all
\(xyz\) axes are \(0\).

The number of ways to choose 2 axes from the 3 axes \(t, R, Q\) is
\(\binom{3}{2} = 3\), and the configurations are classified according to
the combination of 2 non-zero axes as follows.

{\def\LTcaptype{none} % do not increment counter
\begin{longtable}[]{@{}ll@{}}
\toprule\noalign{}
Non-zero 2 axes & Configuration form \\
\midrule\noalign{}
\endhead
\bottomrule\noalign{}
\endlastfoot
\(t, R\) & \((0, 0, 0, \pm, \pm, 0)\) \\
\(t, Q\) & \((0, 0, 0, \pm, 0, Q)\) \\
\(R, Q\) & \((0, 0, 0, 0, \pm, Q)\) \\
\end{longtable}
}

Therefore, the number of types of spin-2 configurations in this model is
\textbf{3}.

\begin{center}\rule{0.5\linewidth}{0.5pt}\end{center}

\subsection{6. Enumeration of Spin-3
Configurations}\label{enumeration-of-spin-3-configurations}

Spin-3 configurations have \(n = 3\) by Definition 2.1, meaning all of
\(t, R, Q\) are non-zero, while all \(xyz\) axes are \(0\).

The number of combinations where all 3 axes \(t, R, Q\) are non-zero is
\(\binom{3}{3} = 1\).

{\def\LTcaptype{none} % do not increment counter
\begin{longtable}[]{@{}ll@{}}
\toprule\noalign{}
Non-zero 3 axes & Configuration form \\
\midrule\noalign{}
\endhead
\bottomrule\noalign{}
\endlastfoot
\(t, R, Q\) & \((0, 0, 0, \pm, \pm, Q)\) (\(Q \neq 0\)) \\
\end{longtable}
}

Therefore, the number of types of spin-3 configurations in this model is
\textbf{1}.

\begin{center}\rule{0.5\linewidth}{0.5pt}\end{center}

\subsection{7. Sign Product of Non-Zero
Axes}\label{sign-product-of-non-zero-axes}

\textbf{Definition 7.1 (Sign product)}. For a boson configuration, when
the \(n\) axes among \(t, R, Q\) with non-zero sign configurations each
take the sign \(+\) or \(-\), the product of these signs is defined as
the \textbf{sign product} \(P(n)\).

\[P(n) = (-1)^n\]

{\def\LTcaptype{none} % do not increment counter
\begin{longtable}[]{@{}cc@{}}
\toprule\noalign{}
\(n\) & \(P(n)\) \\
\midrule\noalign{}
\endhead
\bottomrule\noalign{}
\endlastfoot
0 & \(+1\) \\
1 & \(-1\) \\
2 & \(+1\) \\
3 & \(-1\) \\
\end{longtable}
}

\textbf{Proposition 7.1}. \(P(n)\) depends only on the parity of \(n\).
When \(n\) is even, \(P(n) = +1\); when \(n\) is odd, \(P(n) = -1\).

\textbf{Remark}. \(n = 0\) (no sign reversal) and \(n = 2\) (double
reversal) give the same \(P = +1\), but their algebraic structures
differ. \(n = 0\) is the identity transformation, while \(n = 2\) is the
self-inverse of sign reversal.

The physical interpretation of \(P(n)\) (correspondence to force
directionality) is treated in §10 (an example of interpretation).

\begin{center}\rule{0.5\linewidth}{0.5pt}\end{center}

\subsection{\texorpdfstring{8. Signed Area of Square Faces Containing
the \(R\)
Axis}{8. Signed Area of Square Faces Containing the R Axis}}\label{signed-area-of-square-faces-containing-the-r-axis}

\subsubsection{8.1 Definition of Signed
Area}\label{definition-of-signed-area}

\textbf{Definition 8.1 (Signed area)}. The \textbf{signed area}
\(M(\sigma)\) of a configuration structure \(\sigma\) is defined as the
sum of the products of sign configurations and areas over all square
faces containing the \(R\) axis.

\[M(\sigma) = \sum_{R\text{-containing square face } f} \mathrm{sgn}(f) \cdot \mathrm{Area}(f)\]

\(M(\sigma)\) is a fundamental conserved quantity of this model.
Physical interpretation (such as analogy with mass) is treated in §10.

\subsubsection{\texorpdfstring{8.2 Classification of Faces Containing
the \(R\)
Axis}{8.2 Classification of Faces Containing the R Axis}}\label{classification-of-faces-containing-the-r-axis}

Among the 80 square faces, the faces containing \(R\) (\(= x_5\)) are
classified as follows.

{\def\LTcaptype{none} % do not increment counter
\begin{longtable}[]{@{}llcc@{}}
\toprule\noalign{}
Face type & Bivector & Number of types & Number of faces \\
\midrule\noalign{}
\endhead
\bottomrule\noalign{}
\endlastfoot
Spatial\(-R\) & \(e_i \wedge e_5\) & 3 & \(3 \times 8 = 24\) \\
\(t-R\) & \(e_4 \wedge e_5\) & 1 & 8 \\
\textbf{Total} & & 4 & \textbf{32} \\
\end{longtable}
}

The faces containing the \(R\) axis are 32 out of 80 total faces,
constituting \(2/5\) of the total.

\subsubsection{\texorpdfstring{8.3 Existence of Configurations with
\(M(\sigma) = 0\)}{8.3 Existence of Configurations with M(\textbackslash sigma) = 0}}\label{existence-of-configurations-with-msigma-0}

\textbf{Proposition 8.1}. Among the configuration structures \(\sigma\),
there exist those for which \(M(\sigma) = 0\) (the signed area
completely cancels).

The interpretation of configurations with \(M(\sigma) = 0\) and their
physical correspondence is treated in §10.

\begin{center}\rule{0.5\linewidth}{0.5pt}\end{center}

\subsection{9. Conclusions of the Main
Body}\label{conclusions-of-the-main-body}

In this paper, we defined a 6-dimensional structure obtained by adding a
discrete label axis \(Q\) to the 5-dimensional hyperrectangle, and
derived the following combinatorial properties from its configuration
structure.

\begin{enumerate}
\def\labelenumi{\arabic{enumi}.}
\tightlist
\item
  \textbf{Element counts of hyperrectangles} (§1.2): The number of
  \(j\)-dimensional faces of the \(k\)-dimensional hyperrectangle is
  \(f_j(k) = \binom{k}{j} \cdot 2^{k-j}\). For \(k = 5\): 32 vertices,
  80 edges, 80 square faces, 40 cubic cells, 10 4-facets.
\item
  \textbf{Complete symmetry of 5 axes} (Proposition 1.1): The 5
  geometric axes are all equivalent in combinatorial structure, and any
  axis permutation gives an automorphism of the hyperrectangle.
\item
  \textbf{Axis notation labels} (Definition 1.2): The notation labels
  \(x, y, z, t, R\) are assigned to the 5 axes.
\item
  \textbf{Rigorous definition of sign configurations} (§1.3.3): Range
  constraints on each axis are expressed using the 5 symbols
  \(+, -, 0, \pm, *\).
\item
  \textbf{Definition of spin} (Definition 2.1): Spin \(s\) is classified
  into 8 types (\(s = 0, 1/2, 1, 3/2, 2, 5/2, 3, 7/2\)) from the
  configuration \((x, y, z, t, R, Q)\), based on the number \(n\)
  (\(= 0, 1, 2, 3\)) of non-zero axes among \(t, R, Q\) and the
  structure of the \(xyz\) axes.
\item
  \textbf{Enumeration of \(Q\)-label transitions} (§3): Transitions from
  3 colors to 3 colors give \(3 \times 3 = 9\).
\item
  \textbf{Enumeration of spin-1 configurations} (§4): 2 states with
  \(t\) fixed, 2 states with \(R\) fixed, 9 states with \(Q\) fixed,
  totaling 13 states.
\item
  \textbf{3 types of spin-2 configurations} (§5): \(tR, tQ, RQ\) --- 3
  types.
\item
  \textbf{1 type of spin-3 configurations} (§6): \(tRQ\) --- 1 type.
\item
  \textbf{Sign product \(P(n) = (-1)^n\)} (Definition 7.1, Proposition
  7.1): The product of signs of non-zero axes depends only on the parity
  of \(n\).
\item
  \textbf{Signed area \(M(\sigma)\)} (Definition 8.1): A fundamental
  conserved quantity of the model, defined from the 32 square faces
  containing the \(R\) axis. Configurations with \(M(\sigma) = 0\)
  exist.
\end{enumerate}

All of these are results derived purely from combinatorics and geometry
under the external assumptions of this model (adoption of the
5-dimensional hyperrectangle, addition of the \(Q\) axis).

\begin{center}\rule{0.5\linewidth}{0.5pt}\end{center}

\subsection{10. An Example of Interpretation: Correspondence with the
Standard
Model}\label{an-example-of-interpretation-correspondence-with-the-standard-model}

\textbf{The content of this section is not a claim of this paper but an
example of interpretation. It does not attempt to physically derive the
Standard Model from this model.} This section presents an example of how
the 6-dimensional structure and its mathematical properties defined and
derived in §1--§9 might be mapped to the particle classification of the
Standard Model.

The following correspondences are not justified from within the model
and are \textbf{choices of interpretation from outside}. If an
alternative mapping becomes available in the future, it can be added in
parallel with this section.

\subsubsection{10.1 Assignment of Physical
Labels}\label{assignment-of-physical-labels}

As stated in §1.3.2, the 5 axes are all equivalent in combinatorial
structure, but in the mapping to the Standard Model, the following
labels are assigned.

\begin{itemize}
\tightlist
\item
  \(x_1, x_2, x_3 \to xyz\): 3-dimensional spatial axes
\item
  \(x_4 \to t\): Time axis
\item
  \(x_5 \to R\): Curvature axis (corresponding to the curvature radius
  of space)
\item
  \(Q\): Color label axis (\(0\) = color-neutral, \(r, g, b\) = color
  charge)
\end{itemize}

\subsubsection{10.2 Correspondence to
Quarks/Leptons}\label{correspondence-to-quarksleptons}

The classification of particles corresponds to the value of the \(Q\)
axis as follows.

{\def\LTcaptype{none} % do not increment counter
\begin{longtable}[]{@{}cl@{}}
\toprule\noalign{}
\(Q\) value & Correspondence \\
\midrule\noalign{}
\endhead
\bottomrule\noalign{}
\endlastfoot
\(0\) & Lepton (no color charge) \\
\(r, g, b\) & Quark (color charge present) \\
\end{longtable}
}

For the same geometric configuration \((x, y, z, t, R)\), the existence
of 4 possible values of \(Q\) causes each generation of fermions to
branch into 1 lepton type and 3 quark colors.

\textbf{Charge correspondence}: The sign of electric charge is
distinguished by the sign configuration of the spatial axes. The
denominator 3 of the fractional quark charges (\(\pm 1/3, \pm 2/3\))
coincides with the number of non-zero values of the \(Q\) axis
(\(|\{r, g, b\}| = 3\)). The complete derivation of charge values,
particularly the asymmetry between \(+2/3\) and \(-1/3\), remains an
open problem (§10.10).

\textbf{Consideration of a 7th axis}: The design of adding electric
charge as an independent 7th axis was also considered, but increasing
the number of non-spatial axes to 4 would increase the number of
\(n = 2\) types to \(\binom{4}{2} = 6\), producing a particle spectrum
similar to supersymmetric theories. Given the current experimental
situation (non-discovery of supersymmetric particles), this paper adopts
the approach of reading electric charge from the structure of the \(Q\)
axis and the signs of the spatial axes.

\subsubsection{10.3 Correspondence to 8 Gluon
States}\label{correspondence-to-8-gluon-states}

Of the 9 states of \(Q\)-label transitions enumerated in §3, the
\textbf{8 states excluding the color singlet
(\(r\bar{r} + g\bar{g} + b\bar{b}\))} are mapped to \textbf{gluons
\(g_1, \ldots, g_8\)}.

In the Standard Model, the 8 gluon states are derived as the
8-dimensional adjoint representation of the Lie algebra
\(\mathfrak{su}(3)\) of \(\mathrm{SU}(3)\). In this interpretation, the
same number 8 arises combinatorially from the 4-label structure of the
\(Q\) axis, but the non-commutative structure constants \(f^{abc}\) of
\(\mathrm{SU}(3)\) are not derived from this model.

\subsubsection{10.4 Correspondence to 12 Gauge Boson
States}\label{correspondence-to-12-gauge-boson-states}

Of the 13 spin-1 configuration states enumerated in §4, the 8 states
obtained by excluding the color singlet from the 9 \(Q\)-fixed states,
together with the 2 \(t\)-fixed states and 2 \(R\)-fixed states, give
\textbf{12 states mapped to gauge bosons}.

{\def\LTcaptype{none} % do not increment counter
\begin{longtable}[]{@{}llc@{}}
\toprule\noalign{}
Fixed axis & Correspondence & Number of states \\
\midrule\noalign{}
\endhead
\bottomrule\noalign{}
\endlastfoot
\(t\) & \(W^+, W^-\) & 2 \\
\(R\) & \(\gamma, Z^0\) & 2 \\
\(Q\) & \(g_1, \ldots, g_8\) & 8 \\
\textbf{Total} & & \textbf{12} \\
\end{longtable}
}

The 2 states of the \(W\) boson (\(t = +1, -1\)) correspond to the
positive and negative charges, while \(\gamma\) and \(Z^0\) correspond
to the presence or absence of \(M(\sigma)\) due to the sign of the \(R\)
axis (\(M(\gamma) = 0, M(Z) \neq 0\)). The total of 12 matches the
Standard Model gauge boson state count
\(\gamma + g \times 8 + W^+ + W^- + Z^0 = 12\).

\subsubsection{10.5 Interpretation of Spin-2
Configurations}\label{interpretation-of-spin-2-configurations}

The following correspondences are given to the 3 types of spin-2
configurations enumerated in §5.

{\def\LTcaptype{none} % do not increment counter
\begin{longtable}[]{@{}
  >{\raggedright\arraybackslash}p{(\linewidth - 4\tabcolsep) * \real{0.1765}}
  >{\raggedright\arraybackslash}p{(\linewidth - 4\tabcolsep) * \real{0.3529}}
  >{\raggedright\arraybackslash}p{(\linewidth - 4\tabcolsep) * \real{0.4706}}@{}}
\toprule\noalign{}
\begin{minipage}[b]{\linewidth}\raggedright
Type
\end{minipage} & \begin{minipage}[b]{\linewidth}\raggedright
Fixed axes
\end{minipage} & \begin{minipage}[b]{\linewidth}\raggedright
Correspondence
\end{minipage} \\
\midrule\noalign{}
\endhead
\bottomrule\noalign{}
\endlastfoot
\(tR\) type & \(t, R\) & \textbf{Graviton} (correspondence within the
Standard Model + gravity framework) \\
\(tQ\) type & \(t, Q\) & No known particle in the Standard Model \\
\(RQ\) type & \(R, Q\) & No known particle in the Standard Model \\
\end{longtable}
}

By Proposition 7.1, all 3 types have \(P(2) = +1\), and under the
interpretation of §10.7, they mediate only attractive forces. For the
\(tR\) type, this is consistent with the property of gravity (attraction
only). The \(tQ\) and \(RQ\) types are states whose existence is derived
from the combinatorial structure of this model, but have no counterpart
in the current Standard Model.

\subsubsection{10.6 Interpretation of Spin-3
Configurations}\label{interpretation-of-spin-3-configurations}

The 1 type (\(tRQ\)) of spin-3 configurations enumerated in §6 has no
known particle counterpart in the Standard Model.

\subsubsection{\texorpdfstring{10.7 Correspondence Between the Sign
Product \(P(n)\) and Force
Directionality}{10.7 Correspondence Between the Sign Product P(n) and Force Directionality}}\label{correspondence-between-the-sign-product-pn-and-force-directionality}

The following physical interpretation is given to the sign product
\(P(n) = (-1)^n\) defined in §7.

\begin{itemize}
\tightlist
\item
  \(P(n) = +1\) (\(n\) even): Force is unidirectional
  (\textbf{attraction only})
\item
  \(P(n) = -1\) (\(n\) odd): Force is bidirectional (\textbf{both
  attraction and repulsion})
\end{itemize}

{\def\LTcaptype{none} % do not increment counter
\begin{longtable}[]{@{}ccc@{}}
\toprule\noalign{}
\(n\) & \(P(n)\) & Interpretation of force directionality \\
\midrule\noalign{}
\endhead
\bottomrule\noalign{}
\endlastfoot
0 & \(+1\) & No direction (scalar field, \(s = 0\)) \\
1 & \(-1\) & Both attraction and repulsion (\(s = 1\)) \\
2 & \(+1\) & Attraction only (\(s = 2\)) \\
3 & \(-1\) & Both attraction and repulsion (\(s = 3\)) \\
\end{longtable}
}

Under this interpretation, all 3 types of spin-2 configurations
(\(tR, tQ, RQ\)) have \(n = 2\) (even) and mediate only attractive
forces. This is consistent with the property of gravity (attraction
only). It is also consistent with the electromagnetic force at spin 1
(\(n = 1\), both attraction and repulsion).

\textbf{Remark}. This correspondence is not justified from within the
model, and the physical mechanism connecting the parity of \(P(n)\) to
force directionality is not derived in this paper.

\subsubsection{\texorpdfstring{10.8 Range of the \(R\) Axis and Force
Coupling
Strengths}{10.8 Range of the R Axis and Force Coupling Strengths}}\label{range-of-the-r-axis-and-force-coupling-strengths}

\textbf{Range of the \(R\) axis}: In prior work {[}W5{]}, the
number-theoretic conditions of central projection in discrete space lead
to the discretization \(R = \pm n^2 \cdot L_0\) (\(n \in \mathbb{N}\),
\(L_0\) is the unit length). However, the results of this paper do not
require this discretization and hold even when \(R\) takes continuous
values. Interpreting the range of the \(R\) axis as unbounded
corresponds to the unboundedness of the mass spectrum.

\textbf{Interpretation of each force and the axis correspondence in this
model}:

{\def\LTcaptype{none} % do not increment counter
\begin{longtable}[]{@{}
  >{\raggedright\arraybackslash}p{(\linewidth - 6\tabcolsep) * \real{0.1061}}
  >{\raggedright\arraybackslash}p{(\linewidth - 6\tabcolsep) * \real{0.3030}}
  >{\raggedright\arraybackslash}p{(\linewidth - 6\tabcolsep) * \real{0.2121}}
  >{\raggedright\arraybackslash}p{(\linewidth - 6\tabcolsep) * \real{0.3788}}@{}}
\toprule\noalign{}
\begin{minipage}[b]{\linewidth}\raggedright
Force
\end{minipage} & \begin{minipage}[b]{\linewidth}\raggedright
Mediating particle
\end{minipage} & \begin{minipage}[b]{\linewidth}\raggedright
Involved axis
\end{minipage} & \begin{minipage}[b]{\linewidth}\raggedright
Coupling characteristics
\end{minipage} \\
\midrule\noalign{}
\endhead
\bottomrule\noalign{}
\endlastfoot
Strong force & \(g_1\)--\(g_8\) & \(Q\) axis (8 color transition
channels) & Between color-charged particles only \\
Electromagnetic force & \(\gamma\) & \(R\) axis & Between non-zero
charged particles \\
Weak force & \(W^\pm, Z^0\) & \(t\) axis & All fermions \\
Gravity & \(G\) & \(t \wedge R\) (\(n = 2\)) & All particles
(\(Q\)-independent) \\
\end{longtable}
}

\textbf{Remark}. The coupling constants in the Standard Model near
\(M_Z\) are \(\alpha_s \approx 0.1\), \(\alpha \approx 1/137\),
\(\alpha_W \approx 1/30\), \(\alpha_G \approx 10^{-39}\); the strong,
electromagnetic, and weak forces are of similar order, while gravity
alone is extremely weak. The quantitative derivation of this hierarchy
is an open problem for future work.

\subsubsection{\texorpdfstring{10.9 Analogy Between \(M(\sigma)\) and
Mass}{10.9 Analogy Between M(\textbackslash sigma) and Mass}}\label{analogy-between-msigma-and-mass}

The signed area \(M(\sigma)\) defined in §8.1 has the following
\textbf{analogies} with particle mass.

\begin{itemize}
\tightlist
\item
  Configurations with \(M(\sigma) = 0\) can be interpreted as
  corresponding to massless particles (photon, gluons)
\item
  The unboundedness of the \(R\) axis range corresponds to the
  unboundedness of the mass spectrum
\end{itemize}

However, this paper does not identify \(M(\sigma)\) with mass. The
derivation of specific mass values is an open problem for future work
(§10.12).

\subsubsection{10.10 Table A: Correspondence Between Confirmed Particles
and 6-Dimensional
Configurations}\label{table-a-correspondence-between-confirmed-particles-and-6-dimensional-configurations}

\textbf{Bosons}:

{\def\LTcaptype{none} % do not increment counter
\begin{longtable}[]{@{}
  >{\raggedright\arraybackslash}p{(\linewidth - 10\tabcolsep) * \real{0.2222}}
  >{\raggedright\arraybackslash}p{(\linewidth - 10\tabcolsep) * \real{0.1778}}
  >{\centering\arraybackslash}p{(\linewidth - 10\tabcolsep) * \real{0.1111}}
  >{\raggedright\arraybackslash}p{(\linewidth - 10\tabcolsep) * \real{0.1111}}
  >{\centering\arraybackslash}p{(\linewidth - 10\tabcolsep) * \real{0.1111}}
  >{\centering\arraybackslash}p{(\linewidth - 10\tabcolsep) * \real{0.2667}}@{}}
\toprule\noalign{}
\begin{minipage}[b]{\linewidth}\raggedright
Particle
\end{minipage} & \begin{minipage}[b]{\linewidth}\raggedright
Symbol
\end{minipage} & \begin{minipage}[b]{\linewidth}\centering
\(s\)
\end{minipage} & \begin{minipage}[b]{\linewidth}\raggedright
Configuration \((x, y, z, t, R, Q)\)
\end{minipage} & \begin{minipage}[b]{\linewidth}\centering
\(n\)
\end{minipage} & \begin{minipage}[b]{\linewidth}\centering
Mass (MeV)
\end{minipage} \\
\midrule\noalign{}
\endhead
\bottomrule\noalign{}
\endlastfoot
Higgs & \(H^0\) & 0 & \((\pm, 0, 0, 0, 0, 0)\) etc. & 0 & 125,250 \\
\(W\) boson & \(W^+\) & 1 & \((0, 0, 0, +, 0, 0)\) & 1 & 80,379 \\
\(W\) boson & \(W^-\) & 1 & \((0, 0, 0, -, 0, 0)\) & 1 & 80,379 \\
Photon & \(\gamma\) & 1 & \((0, 0, 0, 0, +, 0)\) & 1 & 0 \\
\(Z\) boson & \(Z^0\) & 1 & \((0, 0, 0, 0, -, 0)\) & 1 & 91,188 \\
Gluon & \(g_1\)--\(g_8\) & 1 & \((0, 0, 0, 0, 0, Q_i \to Q_j)\) & 1 &
0 \\
Graviton & \(G\) & 2 & \((0, 0, 0, \pm, \pm, 0)\) & 2 & 0 \\
\end{longtable}
}

\textbf{Leptons} (\(Q = 0\)):

{\def\LTcaptype{none} % do not increment counter
\begin{longtable}[]{@{}
  >{\raggedright\arraybackslash}p{(\linewidth - 10\tabcolsep) * \real{0.1613}}
  >{\raggedright\arraybackslash}p{(\linewidth - 10\tabcolsep) * \real{0.1290}}
  >{\centering\arraybackslash}p{(\linewidth - 10\tabcolsep) * \real{0.0806}}
  >{\centering\arraybackslash}p{(\linewidth - 10\tabcolsep) * \real{0.1935}}
  >{\raggedright\arraybackslash}p{(\linewidth - 10\tabcolsep) * \real{0.2419}}
  >{\raggedright\arraybackslash}p{(\linewidth - 10\tabcolsep) * \real{0.1935}}@{}}
\toprule\noalign{}
\begin{minipage}[b]{\linewidth}\raggedright
Particle
\end{minipage} & \begin{minipage}[b]{\linewidth}\raggedright
Symbol
\end{minipage} & \begin{minipage}[b]{\linewidth}\centering
\(s\)
\end{minipage} & \begin{minipage}[b]{\linewidth}\centering
Generation
\end{minipage} & \begin{minipage}[b]{\linewidth}\raggedright
Configuration
\end{minipage} & \begin{minipage}[b]{\linewidth}\raggedright
Mass (MeV)
\end{minipage} \\
\midrule\noalign{}
\endhead
\bottomrule\noalign{}
\endlastfoot
Electron & \(e^-\) & 1/2 & 1 & \((+, +, 0, 0, 0, 0)\) & 0.511 \\
Electron neutrino & \(\nu_e\) & 1/2 & 1 & \((+, -, 0, 0, 0, 0)\) &
\(< 10^{-6}\) \\
Muon & \(\mu^-\) & 1/2 & 2 & \((+, 0, +, 0, 0, 0)\) & 105.66 \\
Muon neutrino & \(\nu_\mu\) & 1/2 & 2 & \((+, 0, -, 0, 0, 0)\) &
\(< 0.17\) \\
Tau & \(\tau^-\) & 1/2 & 3 & \((0, +, +, 0, 0, 0)\) & 1,776.9 \\
Tau neutrino & \(\nu_\tau\) & 1/2 & 3 & \((0, +, -, 0, 0, 0)\) &
\(< 15.5\) \\
\end{longtable}
}

Anti-leptons correspond to configurations with the signs of the spatial
axes reversed.

\textbf{Quarks} (\(Q \in \{r, g, b\}\)):

{\def\LTcaptype{none} % do not increment counter
\begin{longtable}[]{@{}llccll@{}}
\toprule\noalign{}
Particle & Symbol & \(s\) & Generation & Configuration & Mass (MeV) \\
\midrule\noalign{}
\endhead
\bottomrule\noalign{}
\endlastfoot
Up & \(u\) & 1/2 & 1 & \((+, +, 0, 0, 0, r/g/b)\) & 2.16 \\
Down & \(d\) & 1/2 & 1 & \((+, -, 0, 0, 0, r/g/b)\) & 4.67 \\
Charm & \(c\) & 1/2 & 2 & \((+, 0, +, 0, 0, r/g/b)\) & 1,270 \\
Strange & \(s\) & 1/2 & 2 & \((+, 0, -, 0, 0, r/g/b)\) & 93.4 \\
Top & \(t\) & 1/2 & 3 & \((0, +, +, 0, 0, r/g/b)\) & 172,760 \\
Bottom & \(b\) & 1/2 & 3 & \((0, +, -, 0, 0, r/g/b)\) & 4,180 \\
\end{longtable}
}

Each quark has 3 color states \(Q = r, g, b\). Anti-quarks correspond to
sign-reversed spatial axes and anti-colors of \(Q\).

\textbf{Summary (states with \(s = 0, 1/2, 1, 2\) mappable to the SM)}:

{\def\LTcaptype{none} % do not increment counter
\begin{longtable}[]{@{}
  >{\raggedright\arraybackslash}p{(\linewidth - 8\tabcolsep) * \real{0.2424}}
  >{\centering\arraybackslash}p{(\linewidth - 8\tabcolsep) * \real{0.0758}}
  >{\centering\arraybackslash}p{(\linewidth - 8\tabcolsep) * \real{0.0758}}
  >{\centering\arraybackslash}p{(\linewidth - 8\tabcolsep) * \real{0.3182}}
  >{\centering\arraybackslash}p{(\linewidth - 8\tabcolsep) * \real{0.2879}}@{}}
\toprule\noalign{}
\begin{minipage}[b]{\linewidth}\raggedright
Classification
\end{minipage} & \begin{minipage}[b]{\linewidth}\centering
\(s\)
\end{minipage} & \begin{minipage}[b]{\linewidth}\centering
\(Q\)
\end{minipage} & \begin{minipage}[b]{\linewidth}\centering
Model-derived count
\end{minipage} & \begin{minipage}[b]{\linewidth}\centering
SM correspondence
\end{minipage} \\
\midrule\noalign{}
\endhead
\bottomrule\noalign{}
\endlastfoot
Higgs & 0 & 0 & 1 & 1 \\
Gauge boson (\(t\) type) & 1 & 0 & 2 & \(W^+, W^-\) \\
Gauge boson (\(R\) type) & 1 & 0 & 2 & \(\gamma, Z^0\) \\
Gluon (\(Q\) type) & 1 & transition & 9 & 8 (excluding 1 color
singlet) \\
Graviton (\(tR\) type) & 2 & 0 & 1 & \(G\) (beyond SM) \\
Lepton & 1/2 & 0 & 6 & 6 \\
Anti-lepton & 1/2 & 0 & 6 & 6 \\
Quark & 1/2 & \(r, g, b\) & 18 & 18 \\
Anti-quark & 1/2 & \(\bar{r}, \bar{g}, \bar{b}\) & 18 & 18 \\
\textbf{Subtotal} & & & \textbf{63} & \textbf{62 (61 + G)} \\
\end{longtable}
}

(Model: 63 states = 62 SM-mapped states + 1 color singlet state). Of the
63 states, 62 are mapped to the Standard Model's 61 states + 1 graviton
state. The remaining 1 state (the color singlet of the \(Q\) transition)
is derived by the model but does not correspond to an independent gauge
boson in the Standard Model.

In addition, configurations with \(s = 3/2, 2, 5/2, 3, 7/2\) (Table B,
§10.11) are derived from the model, but no known particles in the
Standard Model correspond to them.

\subsubsection{10.11 Table B: Configurations Without Standard Model
Counterparts}\label{table-b-configurations-without-standard-model-counterparts}

The following lists states that are derived from the combinatorial
structure of this model but have no known particle counterpart in the
Standard Model. These are combinatorial consequences within the model
and do not assert the existence of undiscovered particles.

\textbf{Bosons (\(xyz\) all \(0\))}:

{\def\LTcaptype{none} % do not increment counter
\begin{longtable}[]{@{}
  >{\raggedright\arraybackslash}p{(\linewidth - 8\tabcolsep) * \real{0.1111}}
  >{\centering\arraybackslash}p{(\linewidth - 8\tabcolsep) * \real{0.0926}}
  >{\centering\arraybackslash}p{(\linewidth - 8\tabcolsep) * \real{0.0926}}
  >{\raggedright\arraybackslash}p{(\linewidth - 8\tabcolsep) * \real{0.3704}}
  >{\centering\arraybackslash}p{(\linewidth - 8\tabcolsep) * \real{0.3333}}@{}}
\toprule\noalign{}
\begin{minipage}[b]{\linewidth}\raggedright
Type
\end{minipage} & \begin{minipage}[b]{\linewidth}\centering
\(s\)
\end{minipage} & \begin{minipage}[b]{\linewidth}\centering
\(n\)
\end{minipage} & \begin{minipage}[b]{\linewidth}\raggedright
Configuration form
\end{minipage} & \begin{minipage}[b]{\linewidth}\centering
Number of states
\end{minipage} \\
\midrule\noalign{}
\endhead
\bottomrule\noalign{}
\endlastfoot
\(tQ\) type & 2 & 2 & \((0, 0, 0, \pm, 0, Q_i \to Q_j)\) &
\(2 \times 9 = 18\) \\
\(RQ\) type & 2 & 2 & \((0, 0, 0, 0, \pm, Q_i \to Q_j)\) &
\(2 \times 9 = 18\) \\
\(tRQ\) type & 3 & 3 & \((0, 0, 0, \pm, \pm, Q_i \to Q_j)\) &
\(4 \times 9 = 36\) \\
\end{longtable}
}

\textbf{Fermions (2 of \(xyz\) non-zero, 1 is \(0\))}:

{\def\LTcaptype{none} % do not increment counter
\begin{longtable}[]{@{}
  >{\raggedright\arraybackslash}p{(\linewidth - 10\tabcolsep) * \real{0.0723}}
  >{\centering\arraybackslash}p{(\linewidth - 10\tabcolsep) * \real{0.0602}}
  >{\centering\arraybackslash}p{(\linewidth - 10\tabcolsep) * \real{0.0602}}
  >{\raggedright\arraybackslash}p{(\linewidth - 10\tabcolsep) * \real{0.3253}}
  >{\centering\arraybackslash}p{(\linewidth - 10\tabcolsep) * \real{0.3373}}
  >{\centering\arraybackslash}p{(\linewidth - 10\tabcolsep) * \real{0.1446}}@{}}
\toprule\noalign{}
\begin{minipage}[b]{\linewidth}\raggedright
Type
\end{minipage} & \begin{minipage}[b]{\linewidth}\centering
\(s\)
\end{minipage} & \begin{minipage}[b]{\linewidth}\centering
\(n\)
\end{minipage} & \begin{minipage}[b]{\linewidth}\raggedright
Non-zero non-spatial axes
\end{minipage} & \begin{minipage}[b]{\linewidth}\centering
Spatial-side pattern count
\end{minipage} & \begin{minipage}[b]{\linewidth}\centering
\(Q\) values
\end{minipage} \\
\midrule\noalign{}
\endhead
\bottomrule\noalign{}
\endlastfoot
\(t\) type & 3/2 & 1 & \(t\) & \(3 \times 4 = 12\) & 4 \\
\(R\) type & 3/2 & 1 & \(R\) & 12 & 4 \\
\(Q\) type & 3/2 & 1 & \(Q\) (non-zero) & 12 & 3 \\
\(tR\) type & 5/2 & 2 & \(t, R\) & 12 & 4 \\
\(tQ\) type & 5/2 & 2 & \(t, Q\) & 12 & 3 \\
\(RQ\) type & 5/2 & 2 & \(R, Q\) & 12 & 3 \\
\(tRQ\) type & 7/2 & 3 & \(t, R, Q\) & 12 & 3 \\
\end{longtable}
}

\textbf{Remark}. The spatial-side pattern count of 12 arises from
\(\binom{3}{1} = 3\) ways to choose the zero axis × \(2^2 = 4\) sign
choices for the 2 non-zero axes. The \(Q\) values are 4 when \(Q = 0\)
is included, and 3 (\(\{r, g, b\}\)) when the \(Q\) axis is required to
be non-zero.

\subsubsection{10.12 Open Problems in This
Interpretation}\label{open-problems-in-this-interpretation}

The following are open problems within the scope of this section's
interpretation.

\textbf{10.12.1 Quantitative Computation of \(M(\sigma)\) and Mass
Ratios}

Explicitly computing the signed area \(M(\sigma)\) for each
configuration structure and comparing with known particle mass ratios is
a means of testing this interpretation.

\begin{itemize}
\tightlist
\item
  \(m_\mu / m_e \approx 206.8\): Mass ratio between generations
\item
  \(m_W / m_Z \approx 0.881\): Relation to the Weinberg angle
\item
  \(m_t / m_e \approx 3.4 \times 10^5\): Maximum mass ratio
\end{itemize}

\textbf{10.12.2 Derivation of Gauge Groups}

A correspondence is suggested between the 4-label structure of the \(Q\)
axis and \(\mathrm{SU}(3)\), the 2 values of the \(t\) axis and
\(\mathrm{SU}(2)\), and the continuous scale of the \(R\) axis and
\(\mathrm{U}(1)\), but the rigorous mechanism for the transition from
discrete structure to continuous gauge groups is unproven. In
particular, the number of gluons 8 is reproduced, but the information of
the non-commutative structure constants \(f^{abc}\) is not derived.

\textbf{10.12.3 CKM Matrix and PMNS Matrix}

How inter-generational mixing (CKM matrix, PMNS matrix) is derived as a
breaking of \(S_3\) symmetry remains unsolved.

\textbf{10.12.4 Higgs Mechanism}

The Higgs boson corresponds to \(n = 0\) (spatial axes only), but how
spontaneous symmetry breaking by the Higgs field is understood within
this framework is an open problem for future work.

\textbf{10.12.5 Complete Derivation of Charge Values}

The sign of electric charge can be read from the sign configurations of
the spatial axes, and the unit of the absolute value of charge (1 vs
1/3) from the \(Q\) axis sector, but the asymmetry between \(+2/3\) and
\(-1/3\) is the most prominent open problem of this interpretation.

\textbf{10.12.6 Additional Predictions of \(s = 2, 3\) Bosons}

If \(tQ\) type and \(RQ\) type spin-2 bosons, and the \(tRQ\) type
spin-3 boson exist, their experimental consequences need to be
clarified.

\subsubsection{10.13 Conclusion of the Interpretation
Example}\label{conclusion-of-the-interpretation-example}

This section presented an example mapping the mathematical properties of
the 6-dimensional structure defined and derived in §1--§9 to the
particle classification of the Standard Model.

\begin{enumerate}
\def\labelenumi{\arabic{enumi}.}
\tightlist
\item
  Of the 8 spin types (\(s = 0, 1/2, 1, 3/2, 2, 5/2, 3, 7/2\)), 62 of
  the 63 states with \(s = 0, 1/2, 1, 2\) can be mapped to known
  particles of the Standard Model (61 + graviton) (the remaining 1 state
  is the color singlet)
\item
  The 4 values of the \(Q\) axis correspond to color charge and the
  lepton/quark branching
\item
  Of the 9 \(Q\)-label transitions, 8 states correspond to gluons (the
  color singlet has no counterpart)
\item
  Of the 13 spin-1 configuration states, 12 correspond to Standard Model
  gauge bosons
\item
  The \(tR\) type of spin-2 configurations corresponds to the graviton.
  The \(tQ\) and \(RQ\) types are derived from the model but have no
  Standard Model counterpart
\item
  The \(tRQ\) type of spin-3 configurations similarly has no counterpart
\item
  Fermion types with \(s = 3/2, 5/2, 7/2\) also exist combinatorially
  but do not correspond to known particles (Table B)
\item
  \(M(\sigma)\) has analogies with particle mass
\item
  The types of axes involved in each force can be interpreted as the
  qualitative origin of the coupling strength hierarchy
\end{enumerate}

This model derives configurations that \textbf{exceed} the known
particles of the Standard Model. These ``surplus states'' are open
problems in this interpretation and will be examined in future research
as one of the following: (a) predictions of undiscovered particles, (b)
combinations not physically realized (existence of selection rules), or
(c) limitations of this interpretation.

The derivation of specific mass values, charge values, non-commutative
structure of gauge groups, and generation mixing matrices remains an
open problem for future work (§10.12).

\begin{center}\rule{0.5\linewidth}{0.5pt}\end{center}

\subsection{References}\label{references}

{[}1{]} Finkelstein, D. \& Rubinstein, J. ``Connection between Spin,
Statistics, and Kinks.'' \emph{J. Math. Phys.} 9, 1762 (1968).

{[}2{]} Atiyah, M.F., Manton, N.S. \& Schroers, B.J. ``Geometric Models
of Matter.'' \emph{Proc. R. Soc. A} 468, 1252 (2012). arXiv: 1108.5151.

{[}3{]} Furey, C. ``Standard Model Physics from an Algebra?'' arXiv:
1611.09182 (2016).

{[}4{]} Kaluza, Th. ``Zum Unitätsproblem der Physik.''
\emph{Sitzungsber. Preuss. Akad. Wiss. Berlin} 966--972 (1921).

{[}5{]} Coxeter, H.S.M. \emph{Regular Polytopes}. Dover (1973).

{[}6{]} Kihara, N. ``Bivectorial Classification of Spin Derived from the
Orientation Structure of the 5-Dimensional Hypercube.'' DOI:
\href{https://doi.org/10.5281/zenodo.19643358}{10.5281/zenodo.19643358}
(2026).

{[}7{]} Kihara, N. ``Geometric Classification of Spin Derived from the
Orientation Structure of the 5-Dimensional Orthoplex.'' DOI:
\href{https://doi.org/10.5281/zenodo.19630972}{10.5281/zenodo.19630972}
(2026).

{[}W5{]} Kihara, N. ``Complete Spherical Coverage of Central Projection
in Discrete Space and the Number-Theoretic Necessity of the
5-Dimensional Background Space.'' DOI:
\href{https://doi.org/10.5281/zenodo.19624957}{10.5281/zenodo.19624957}
(2026).

\end{document}
